% !TeX program = xelatex
\documentclass[UTF8,a4paper,zihao=-4,fontset=none]{ctexart}
\setCJKmainfont[BoldFont={LXGW Neo XiHei}]{LXGW Neo ZhiSong}
\setCJKsansfont[AutoFakeBold=2]{LXGW Neo XiHei}
\setCJKmonofont{LXGW Neo XiHei}
\setCJKfamilyfont{zhhei}[AutoFakeBold=2]{LXGW Neo XiHei}
\providecommand{\heiti}{\CJKfamily{zhhei}}
% Shared renderer entrypoint for generated lessonplan.tex files.
% Generated documents should load ctexart, then input this file.

\usepackage{hepingjie_lessonplan}
% Hepingjie lesson-plan overlay coordinates.
% Coordinate origin: top-left corner of an A4 portrait page.
% Units: mm. These coordinates are frozen for templates/lessonplan/hepingjie_blank.pdf.

\newcommand{\LPBackgroundPdf}{templates/lessonplan/hepingjie_blank.pdf}

% Page 1 fixed fields.
\newcommand{\LPDateYearX}{65mm}
\newcommand{\LPDateMonthX}{89mm}
\newcommand{\LPDateDayX}{110mm}
\newcommand{\LPDateY}{38mm}
\newcommand{\LPRecordPageX}{170mm}
\newcommand{\LPRecordPageY}{38mm}

\newcommand{\LPTopicX}{32mm}
\newcommand{\LPTopicY}{49mm}
\newcommand{\LPTopicW}{108mm}
\newcommand{\LPTypeX}{154mm}
\newcommand{\LPTypeY}{49mm}
\newcommand{\LPTypeW}{40mm}

\newcommand{\LPUnitTotalX}{67mm}
\newcommand{\LPUnitTotalY}{61mm}
\newcommand{\LPTopicTotalX}{111mm}
\newcommand{\LPTopicTotalY}{61mm}
\newcommand{\LPCurrentPeriodX}{169mm}
\newcommand{\LPCurrentPeriodY}{61mm}

\newcommand{\LPBodyX}{33mm}
\newcommand{\LPObjectiveY}{74mm}
\newcommand{\LPObjectiveH}{39mm}
\newcommand{\LPKeyPointY}{119mm}
\newcommand{\LPKeyPointH}{15mm}
\newcommand{\LPDifficultyY}{137mm}
\newcommand{\LPDifficultyH}{15mm}
\newcommand{\LPMethodY}{156mm}
\newcommand{\LPMethodH}{15mm}
\newcommand{\LPMediaY}{174mm}
\newcommand{\LPMediaH}{15mm}
\newcommand{\LPBlackboardY}{196mm}
\newcommand{\LPBlackboardH}{73mm}
\newcommand{\LPBodyW}{164mm}

% Teaching process boxes on pages 2-4.
\newcommand{\LPProcessX}{33mm}
\newcommand{\LPProcessY}{31mm}
\newcommand{\LPProcessW}{163mm}
\newcommand{\LPProcessFullH}{249mm}
\newcommand{\LPProcessPageFourH}{160mm}


\usetikzlibrary{calc}
\begin{document}

\LPBackgroundPage{1}{
  \LPCenterBox{\LPDateYearX}{\LPDateY}{12mm}{5mm}{}
  \LPCenterBox{\LPDateMonthX}{\LPDateY}{10mm}{5mm}{}
  \LPCenterBox{\LPDateDayX}{\LPDateY}{10mm}{5mm}{}
  \LPCenterBox{\LPRecordPageX}{\LPRecordPageY}{10mm}{5mm}{1}
  \LPHeaderCell{\LPTopicX}{\LPTopicY}{\LPTopicW}{8mm}{移项解一元一次方程}
  \LPHeaderCell{\LPTypeX}{\LPTypeY}{\LPTypeW}{8mm}{运算技能课}
  \LPCenterBox{\LPUnitTotalX}{\LPUnitTotalY}{12mm}{5mm}{8}
  \LPCenterBox{\LPTopicTotalX}{\LPTopicTotalY}{12mm}{5mm}{1}
  \LPCenterBox{\LPCurrentPeriodX}{\LPCurrentPeriodY}{12mm}{5mm}{1}
  \LPTextBox{\LPBodyX}{\LPObjectiveY}{\LPBodyW}{\LPObjectiveH}{\LPPlainLine{说明移项来源于等式的性质，解释“移项要变号”的理由。}\LPPlainLine{按移项、合并同类项、系数化为 1 的步骤解一元一次方程。}\LPPlainLine{判断并订正移项不变号、负系数处理错误，能用代入法检验方程的解。}\LPPlainLine{从简单情境中找出相等关系并列一元一次方程。}}
  \LPTextBox{\LPBodyX}{\LPKeyPointY}{\LPBodyW}{\LPKeyPointH}{从等式性质理解移项，并按规范步骤解一元一次方程。}
  \LPTextBox{\LPBodyX}{\LPDifficultyY}{\LPBodyW}{\LPDifficultyH}{准确识别移动的项及其原符号，解释变号依据并正确处理负系数。}
  \LPTextBox{\LPBodyX}{\LPMethodY}{\LPBodyW}{\LPMethodH}{复习迁移；对比推理；板演纠错}
  \LPTextBox{\LPBodyX}{\LPMediaY}{\LPBodyW}{\LPMediaH}{PPT；黑板；步骤卡；反馈卡}
  \LPTextBox{\LPBodyX}{\LPBlackboardY}{\LPBodyW}{\LPBlackboardH}{\LPBoardTitle{移项解一元一次方程}\LPBoardLine{1}{依据：等式两边同时加或减同一个式子。}\LPBoardLine{2}{移项：把等式一边的某项变号后移到另一边。}\LPBoardLine{3}{步骤：移项 → 合并同类项 → 系数化为 1 → 检验。}\LPBoardLine{4}{例：$3x+7=32-2x$，移项得 $3x+2x=32-7$。}\LPBoardLine{5}{易错：移动的项不变号；$-ax=b$ 除以负数时丢负号。}}
}

\LPBackgroundPage{2}{
  \LPProcessBox{\LPProcessX}{\LPProcessY}{\LPProcessW}{\LPProcessFullH}{\LPProcessRow{复习等式性质}{\LPTeacherPart{问题}{由 $x+5=12$ 得 $x=7$，等式两边同时做了什么？；要消去 $3x+7=32-2x$ 右边的 $-2x$，两边应同时加什么？}\LPTeacherPart{组织}{让学生口头说明每一步依据，把“同时做同一种运算”板书在等号两侧。}\LPTeacherPart{说明}{等式两边同时加或减同一个式子，结果仍相等。}\LPTeacherPart{对应练习}{口答：要消去左边的 $+7$，等式两边应同时减去 7。}\LPTeacherPart{纠错}{若只说“把 5 挪过去”，追问“等式两边实际做了什么”。}\LPTeacherPart{归纳}{方程变形必须保持等式成立。}\LPTeacherPart{意图}{激活等式性质，为解释移项变号提供依据。}}{\LPStudentPart{活动}{解热身方程并完整说出等式性质。；在原方程两边写出相同的加减式。}\LPStudentPart{预期}{两边同时减 5；要消去 $-2x$，两边同时加 $2x$。}}{5}
\LPProcessRow{从等式性质发现移项}{\LPTeacherPart{问题}{比较完整变形与简写，$-2x$ 和 $+7$ 的位置、符号发生了什么变化？；符号变化是凭空规定，还是省略了哪一步？}\LPTeacherPart{组织}{先完整写出两边同时加 $2x$、再同时减 7 的过程，再与 $3x+2x=32-7$ 对照。}\LPTeacherPart{说明}{把等式一边的某项变号后移到另一边，叫作移项；它是等式性质变形的简写。}\LPTeacherPart{例题}{$3x+7=32-2x \Rightarrow 3x+2x=32-7$。}\LPTeacherPart{对应练习}{将 $5x-4=2x+8$ 只做移项，写成 $5x-2x=8+4$，并说明两次变号。}\LPTeacherPart{纠错}{强调“项”包含前面的符号；没有移动的项不能随意变号。}\LPTeacherPart{归纳}{移项改变位置，同时改变该项的符号。}\LPTeacherPart{意图}{通过两种写法对比，让“移项变号”有可说明的数学依据。}}{\LPStudentPart{活动}{圈出移动的两项，比较移项前后的符号。；先写完整等式性质，再写移项简写。}\LPStudentPart{预期}{$-2x$ 移到左边变为 $+2x$，$+7$ 移到右边变为 $-7$；本质是等式两边同时加减同一个式子。}}{8}
\LPProcessRow{例题示范与对应练习}{\LPTeacherPart{问题}{解 $3x+7=32-2x$ 时，哪些项应放在同一边？；每一步的任务分别是什么？}\LPTeacherPart{组织}{按“移项—合并同类项—系数化为 1—检验”逐行板演，并让学生口述依据。}\LPTeacherPart{说明}{$3x+2x=32-7$，$5x=25$，$x=5$；代入得两边都等于 22。}\LPTeacherPart{例题}{解：$\begin{aligned}3x+7&=32-2x,\\3x+2x&=32-7,\\5x&=25,\\x&=5.\end{aligned}$}\LPTeacherPart{对应练习}{解 $4x-3=2x+9$，答案 $x=6$；完成后代入检验。}\LPTeacherPart{纠错}{若学生一步写出答案，要求补写三步并在每步右侧标注名称。}\LPTeacherPart{归纳}{每一步只完成一个明确任务，书写顺序固定。}\LPTeacherPart{意图}{用一例一练形成规范书写流程，并让学生把步骤与依据对应起来。}}{\LPStudentPart{活动}{先说移动哪一项及新符号，再独立完成对应练习。；同桌互查步骤名称和检验结果。}\LPStudentPart{预期}{$4x-2x=9+3$，$2x=12$，$x=6$；代入左右两边都为 21。}}{7}}
}

\LPBackgroundPage{3}{
  \LPProcessBox{\LPProcessX}{\LPProcessY}{\LPProcessW}{\LPProcessFullH}{\LPProcessRow{变式与负系数}{\LPTeacherPart{问题}{解 $x-3=1.5x+1$，移项后为什么得到 $-0.5x=4$？；$-0.5x=4$ 的两边应同除以什么？}\LPTeacherPart{组织}{让学生先独立写前两步，再选取“负号丢失”的板演进行比较。}\LPTeacherPart{说明}{$x-1.5x=1+3$，$-0.5x=4$，两边同除以 $-0.5$，得 $x=-8$。}\LPTeacherPart{例题}{解：$\begin{aligned}x-3&=1.5x+1,\\x-1.5x&=1+3,\\-0.5x&=4,\\x&=-8.\end{aligned}$}\LPTeacherPart{对应练习}{解 $2x+1=5x-8$，答案 $x=3$；说明负系数出现在哪一步。}\LPTeacherPart{纠错}{把“系数化为 1”说完整：两边同除以未知数的系数，包括负号。}\LPTeacherPart{归纳}{负系数化为 1 时，除数必须连同负号。}\LPTeacherPart{意图}{在紧跟变式中暴露负系数错误，巩固步骤而不是记答案。}}{\LPStudentPart{活动}{独立完成变式，圈出 $-0.5$ 和除数。；代入原方程检验 $x=-8$。}\LPStudentPart{预期}{两边同除以 $-0.5$，所以解为 $-8$；代入两边都为 $-11$。}}{7}
\LPProcessRow{错例辨析与检验}{\LPTeacherPart{问题}{错解“由 $2x-4=12$ 得 $2x=8$”的第一处错误在哪里？；怎样用代入检验快速发现错误？}\LPTeacherPart{组织}{要求学生先判断、再说明理由、最后给出规范订正。}\LPTeacherPart{说明}{$-4$ 移到右边应变为 $+4$，故 $2x=16$，$x=8$；代入检验成立。}\LPTeacherPart{例题}{判断 $5x+6=2x$ 移项为 $5x-2x=-6$：正确，两项移项后均变号。}\LPTeacherPart{对应练习}{判断由 $4x+7=2x-5$ 得 $4x-2x=-5-7$，所以 $x=-6$：正确，并说明两次移项均变号。}\LPTeacherPart{纠错}{不接受只写“错误”，必须指出移动的项、符号和正确等式。}\LPTeacherPart{归纳}{纠错三步：找第一处错误、写依据、代入检验。}\LPTeacherPart{意图}{用错例迫使学生解释规则，并建立代入检验习惯。}}{\LPStudentPart{活动}{标出第一处错误，写出完整订正并代入检验。；同桌互评理由是否包含“哪一项怎样变号”。}\LPStudentPart{预期}{正确订正为 $2x=12+4$，$x=8$；判断题中的两次移项都改变了符号。}}{5}
\LPProcessRow{简单应用}{\LPTeacherPart{问题}{两组共整理 55 本书，甲组比乙组的 2 倍少 5 本，若乙组为 $x$ 本，甲组怎样表示？；哪个相等关系可以列方程？}\LPTeacherPart{组织}{按“设未知数—表示相关量—找相等关系—列方程—作答”组织口头建模。}\LPTeacherPart{说明}{甲组为 $2x-5$，方程 $x+2x-5=55$，解得 $x=20$，甲组 35 本。}\LPTeacherPart{纠错}{若把“少 5”写成 $2x+5$，用数量大小反问检验表达式。}\LPTeacherPart{归纳}{表示同一个总量的关系可以转化为方程。}\LPTeacherPart{意图}{用简短应用题连接教材的问题背景，保持本节重点仍在移项。}}{\LPStudentPart{活动}{写出设句和数量关系，口头完成解方程与答语。}\LPStudentPart{预期}{$x+(2x-5)=55$，乙组 20 本，甲组 35 本。}}{3}}
}

\LPBackgroundPage{4}{
  \LPProcessBox{\LPProcessX}{\LPProcessY}{\LPProcessW}{\LPProcessPageFourH}{\LPProcessRow{当堂检测}{\LPTeacherPart{问题}{哪一步最容易出现符号错误，如何用检验发现？}\LPTeacherPart{组织}{投放 4 题：①说明移项变号依据；②解 $6x+5=2x+17$；③解 $0.5x-2=x+1$；④订正 $3x-7=8\Rightarrow3x=1$。}\LPTeacherPart{说明}{答案：等式性质；$x=3$；$x=-6$；应为 $3x=15$，$x=5$。}\LPTeacherPart{纠错}{同桌互判后举牌统计；低于 3 题正确者重写一道完整“等式性质版本”。}\LPTeacherPart{归纳}{达标标准：4 题正确 3 题及以上，并能说明至少一次变号依据。}\LPTeacherPart{意图}{同时检测概念依据、基本技能、负系数和错例订正。}}{\LPStudentPart{活动}{独立完成，互判后对错误题做代入检验。}\LPStudentPart{预期}{能按三步规范解题并用等式性质解释移项。}}{6}
\LPProcessRow{回顾总结与作业}{\LPTeacherPart{问题}{移项与等式性质是什么关系？；解方程的三步流程是什么？；负系数和移项时各要提醒自己什么？}\LPTeacherPart{组织}{根据学生回答补全板书流程，布置分层作业。}\LPTeacherPart{说明}{基础巩固：解 $2x+3=11$、$5x-4=2x+8$，订正 $x-6=9\Rightarrow x=3$。提升思考：用一段话解释移项为什么变号，并解 $0.2x+1=0.5x-5$。}\LPTeacherPart{纠错}{若学生只背“移项变号”，追问“等式两边实际做了什么”。}\LPTeacherPart{归纳}{移项是等式性质的简写；规范步骤和代入检验共同保证正确。}\LPTeacherPart{意图}{从口诀回到依据，形成可迁移的解方程流程。}}{\LPStudentPart{活动}{用一句话概括移项本质，口述三步流程和一个易错点。；记录分层作业。}\LPStudentPart{预期}{移项相当于等式两边同时加减同一式子；步骤为移项、合并同类项、系数化为 1。}}{4}}
}
\end{document}
